\begin{textsample}{NEG Dim 8 – 1960 – Score -3.00 – 1960/tv\_com\_1960\_86.txt}  \label{ex:f8_neg_006}
At first, a small black-and-silver pointed tool lies diagonally on a glossy, pale blue-white studio surface against a dark background.
It looks like a compact metal punch or probe, with a shiny pointed end and a darker handle.
A light-skinned adult hand reaches in, picks it up, and holds it upright so the bright metal point is emphasized.
Another finger hovers above and then presses down toward the sharp conical tip, making the “ point ” visually literal.
The finger pad fills much of the image as it comes close to the tiny metal point, then lifts away, leaving the polished tip standing out against the dark background.
A one-gallon can of Zerex appears on the same clean studio surface.
The can is white with large black letters spelling “ ZEREX, ” a DuPont oval logo near the top, and the words “ GUARANTEED anti-leak ” above the brand name, with “ anti-freeze \& summer coolant ” below.
The word “ anti-leak ” is printed in orange-red, while most other lettering is black.
A dark, wet-looking vertical streak or glossy strip runs down the right side of the can, suggesting liquid or a leak motif.
The product label fills more and more of the image until the oversized black letters dominate.
The hand returns with the pointed tool and brings its tip toward the front of the can, \textbf{aiming} at the large black letters of “ ZEREX. ” The tool is pressed into the can ’s face near the middle of the label.
A thin, bright stream of liquid suddenly shoots from the puncture toward the left, a narrow white line crossing the label area.
The stream continues for a moment, then weakens and disappears.
A small dark wet spot remains at the puncture site, with only a bead or tiny mark visible.
The hand briefly reaches in with fingers close to the hole, as if checking or wiping the spot, and then the can face is shown with the leak no longer spraying.
The \textbf{demonstration} expands into a more technical-looking setup.
On the left, the Zerex can remains partly visible, cropped so only the left portion of the label and the big “ REX ” letters show.
On the right, white block text appears over the dark blue background : “ 50 \% ZEREX SOLUTION CIRCULATING UNDER PRESSURE. ” A light-skinned hand brings the pointed metal tool in from the right and presses it against the side of the can.
At first, a small wet burst appears where the point touches.
Then the hand jabs again, and liquid sprays outward in bright, white, needle-like jets.
Several streams arc to the right, some straight and forceful, others splashing into fine streaks and droplets.
The hand keeps moving the pointed tool along the can ’s side, making or touching several leak points.
Each puncture produces a sudden spray, with water-like lines shooting across the blue background.
As the tool is withdrawn, multiple arcing streams remain visible, then some thin out, break into streaks, or fade.
The can ’s left edge stays steady while the sprays flare and diminish to the right, visually demonstrating pressurized leaks being challenged by the Zerex solution.
The commercial then returns to the product label, now large and centered again : DuPont at the top, “ GUARANTEED anti-leak, ” “ ZEREX, ” and “ anti-freeze \& summer coolant. ” The composition widens to show the full metal can on the right and a larger white plastic jug on the left, both labeled Zerex anti-leak anti-freeze.
The jug has a molded handle near the top and a cap, while the can is \textbf{cylindrical} with a metal rim.
The same black-and-silver pointed tool lies horizontally on the glossy surface in front of the packages, tying the leak-puncture \textbf{demonstration} back to the product display.
The final image holds the two Zerex containers side by side against the softly lit blue studio background, with the brand name clearly visible on both.

% matched lemmas: aim, cylindrical, demonstration
\end{textsample}
